\clearpage\NoSourceHead\SourceLeaf{10}
\begin{center}\large\bfseries Einleitung.\end{center}
\noindent\hspace*{\parindent}Wer eine Darstellung alles dessen unternimmt, was die mathematischen Wissenschaften \textsc{Gauss} verdanken, muss füglich seine zahlentheoretischen Arbeiten an erster Stelle in Betracht ziehen. Nicht so sehr aus dem Grunde, weil die Disquisitiones arithmeticae, sein grösstes und zwar nicht dem Erscheinen, aber der Entstehung nach erstes Werk, der Zahlentheorie gewidmet sind, als vielmehr deswegen, weil in der Tat in dem Kranze von \textsc{Gauss}’ epochemachenden wissenschaftlichen Leistungen seine Entdeckungen im Gebiete der Zahlentheorie die grossartigsten und in ihrer Wirkung auf die weitere Entwicklung der Mathematik die bedeutendsten gewesen sind; von jenem monumentalen Jugendwerke an datieren wir erst in strengerem Sinne die \textls[100]{Wissenschaft der höheren Arithmetik}.

Im Folgenden wollen wir versuchen, eine zwar möglichst knappe, aber doch das Wesentliche erschöpfende systematische Skizze dieser zahlentheoretischen \textsc{Gauss}schen Ergebnisse zu zeichnen, indem wir zunächst von den hauptsächlichsten Teilen der D. A.\footnote{Zur Abkürzung steht im Folgenden:\par
D. A. für Disquisitiones arithmeticae, A. R. für Analysis Residuorum, T. für Tagebuch von \textsc{Gauss}, W. für \textsc{Gauss}’ Werke.} eine Analyse geben, und dann zeigen, wie die hier begründeten Theorien in den späteren Arbeiten von \textsc{Gauss} sich entwickelt und neue Früchte gezeitigt haben. Zugleich werden wir bemüht sein, dem Interesse an der Entstehungsgeschichte dieser Arbeiten oder einzelner Lehrsätze an der Hand der eigenen \textsc{Gauss}schen Aufzeichnungen darüber, soweit möglich, Rechnung zu tragen.
\HistoricalFoot{}{1*}
\clearpage
